\section{Conclusion}
\label{sec:conclusion}

Even though it is well-known that the OBD depends on the choice of reference group, we show that this dependence can be highly consequential, leading to sign reversals in both the explained and unexplained components.
We present real and simulated examples in which such a reversal occurs.%We present a simple healthcare example, inspired by real-world data, as well as a real clinical case study, in which one OBD direction suggests that quality of care improves patient outcomes, while the other suggests it worsens them. 
We further prove that such sign reversals are not isolated phenomena, but can occur in a wide variety of scenarios. 
However, we demonstrate that while the set of regression parameters $(\beta_g, \alpha_g, \mu_g), g \in \{K, H\}$ leading to sign flips is substrantial, there is reason to believe sign flips will be rare in practice.
We therefore believe that practitioners should be cautiously optimistic about use of the OBD: it is likely to provide robust and useful insights about differences between populations, but it is critical that practitioners compute and report its robustness to the choice of reference group, less they be mislead by an arbitrary choice of the reference group.
%Practitioners should therefore be cautious optimistic when using the OBD, at least in settings where no reference group is clearly justified.
